\documentclass[10pt,a4paper]{amsart}
\usepackage{amsmath,amssymb,amsthm,graphicx}
\usepackage{verbatim}
\newenvironment{reference}{\comment}{\endcomment}
%\newenvironment{reference}{}{}
\newenvironment{slogan}{\comment}{\endcomment}
\newenvironment{history}{\comment}{\endcomment}

% For commutative diagrams we use Xy-pic
\usepackage[all]{xy}

% We use 2cell for 2-commutative diagrams.
\xyoption{2cell}
% Theorem environments.
%
\theoremstyle{plain}
\newtheorem{theorem}[subsection]{Theorem}
\newtheorem{proposition}[subsection]{Proposition}
\newtheorem{lemma}[subsection]{Lemma}

\theoremstyle{definition}
\newtheorem{definition}[subsection]{Definition}
\newtheorem{example}[subsection]{Example}
\newtheorem{exercise}[subsection]{Exercise}
\newtheorem{situation}[subsection]{Situation}

\theoremstyle{remark}
\newtheorem{remark}[subsection]{Remark}
\newtheorem{remarks}[subsection]{Remarks}

\numberwithin{equation}{subsection}

% Macros
%
\def\lim{\mathop{\mathrm{lim}}\nolimits}
\def\colim{\mathop{\mathrm{colim}}\nolimits}
\def\Spec{\mathop{\mathrm{Spec}}}
\def\Hom{\mathop{\mathrm{Hom}}\nolimits}
\def\Ext{\mathop{\mathrm{Ext}}\nolimits}
\def\SheafHom{\mathop{\mathcal{H}\!\mathit{om}}\nolimits}
\def\SheafExt{\mathop{\mathcal{E}\!\mathit{xt}}\nolimits}
\def\Sch{\mathit{Sch}}
\def\Mor{\mathop{\mathrm{Mor}}\nolimits}
\def\Ob{\mathop{\mathrm{Ob}}\nolimits}
\def\Sh{\mathop{\mathit{Sh}}\nolimits}
\def\NL{\mathop{N\!L}\nolimits}
\def\CH{\mathop{\mathrm{CH}}\nolimits}
\def\proetale{{pro\text{-}\acute{e}tale}}
\def\etale{{\acute{e}tale}}
\def\QCoh{\mathit{QCoh}}
\def\Ker{\mathop{\mathrm{Ker}}}
\def\Im{\mathop{\mathrm{Im}}}
\def\Coker{\mathop{\mathrm{Coker}}}
\def\Coim{\mathop{\mathrm{Coim}}}

% Boxtimes
%
\DeclareMathSymbol{\boxtimes}{\mathbin}{AMSa}{"02}

%
% Macros for moduli stacks/spaces
%
\def\QCohstack{\mathcal{QC}\!\mathit{oh}}
\def\Cohstack{\mathcal{C}\!\mathit{oh}}
\def\Spacesstack{\mathcal{S}\!\mathit{paces}}
\def\Quotfunctor{\mathrm{Quot}}
\def\Hilbfunctor{\mathrm{Hilb}}
\def\Curvesstack{\mathcal{C}\!\mathit{urves}}
\def\Polarizedstack{\mathcal{P}\!\mathit{olarized}}
\def\Complexesstack{\mathcal{C}\!\mathit{omplexes}}
% \Pic is the operator that assigns to X its picard group, usage \Pic(X)
% \Picardstack_{X/B} denotes the Picard stack of X over B
% \Picardfunctor_{X/B} denotes the Picard functor of X over B
\def\Pic{\mathop{\mathrm{Pic}}\nolimits}
\def\Picardstack{\mathcal{P}\!\mathit{ic}}
\def\Picardfunctor{\mathrm{Pic}}
\def\Deformationcategory{\mathcal{D}\!\mathit{ef}}
\usepackage{cleveref}
\crefname{originalsection}{original section}{original sections}
\Crefname{originalsection}{Original section}{Original sections}
\makeatletter
\newlabel{section-localize}{{00XZ}{source}{}{sourcepointer-1}{}}
\newlabel{section-localize@cref}{{[originalsection][0][]00XZ}{[1][0][]source}{}{}{}}
\makeatother
\crefname{originallemma}{original lemma}{original lemmas}
\Crefname{originallemma}{Original lemma}{Original lemmas}
\makeatletter
\newlabel{lemma-topos-good-site}{{03CI}{source}{}{sourcepointer-2}{}}
\newlabel{lemma-topos-good-site@cref}{{[originallemma][0][]03CI}{[1][0][]source}{}{}{}}
\makeatother
\crefname{originallemma}{original lemma}{original lemmas}
\Crefname{originallemma}{Original lemma}{Original lemmas}
\makeatletter
\newlabel{lemma-essential-image-j-shriek}{{00Y1}{source}{}{sourcepointer-3}{}}
\newlabel{lemma-essential-image-j-shriek@cref}{{[originallemma][0][]00Y1}{[1][0][]source}{}{}{}}
\makeatother
\crefname{originallemma}{original lemma}{original lemmas}
\Crefname{originallemma}{Original lemma}{Original lemmas}
\makeatletter
\newlabel{lemma-compute-j-shriek-restrict}{{03EE}{source}{}{sourcepointer-4}{}}
\newlabel{lemma-compute-j-shriek-restrict@cref}{{[originallemma][0][]03EE}{[1][0][]source}{}{}{}}
\makeatother
\title{The slice of a topos over a sheaf is a topos with a canonical localization}
\author{The Stacks Project Authors}
\date{stacks-project:sites.tex:local-snapshot-20261012; GNU FDL 1.2 or later}
\begin{document}
\maketitle
\noindent\textit{Editorial note.} All delivered mathematical text is the author's own, copied from the pinned Stacks Project chapter file sites-cohomology.tex; no statement or proof step was rewritten, reordered or removed.\par\medskip
\noindent\textit{Editorial note.} The proof is the author's complete argument; every step the author writes out is inside the counted span and no closing sentence is a placeholder.\par\medskip
\noindent\textit{Editorial note.} Every labelled reference occurring inside the counted spans is carried in assembly.external\_references as a non-load-bearing navigation pointer (same-chapter ones at their definition site, cross-chapter ones at the invocation inside the counted span), so no load-bearing mathematics is replaced by a stub. No bibliography citation occurs inside the counted spans.\par\medskip
\noindent\textit{Editorial note.} The checked chapter list for same-name collisions is the set of stacks chapters staged in this repository: chow, cohomology, divisors, etale-cohomology, flat, more-algebra, more-morphisms, morphisms, spaces-more-morphisms, varieties together with this card's chapter sites-cohomology.\par\medskip
\noindent\textit{Editorial note.} The preamble is the chapter's own retained minimal preamble (sources/stacks-sites-cohomology/retained-preamble.tex), a verbatim excerpt of the native preamble: the Xy-pic setup and the native theorem declarations with the mathematical macros.\par\medskip
\noindent\textit{Allowed standard prerequisites (editorial).}\begin{itemize}
\item Good sites. Same chapter, Lemma lemma-topos-good-site. Reduces the statement to a site with subcanonical topology and a representable sheaf.
\item The equivalence between localized sheaves and slices. Same chapter, Lemma lemma-essential-image-j-shriek. Supplies the identification of the slice with the category of sheaves on the localized site.
\item The computation of the lower shriek of a restriction. Same chapter, Lemma lemma-compute-j-shriek-restrict. Supplies the description of the pullback of the localization.
\end{itemize}\medskip
\ifdefined\hypertarget\hypertarget{sourcepointer-1}{}\fi
\noindent\textit{Source reference.} Original section 00XZ: sources/stacks-sites/pinned-original-source.tex, source lines 7554--7554. Navigation pointer invoked at proof line 7554; the counted span contains the pointer only and the cited result is not reproduced inside it, so the pointer is registered as a non-load-bearing navigation stub in the established form for Stacks material. No bibliography citation is involved.\par\medskip
\ifdefined\hypertarget\hypertarget{sourcepointer-2}{}\fi
\noindent\textit{Source reference.} Original lemma 03CI: sources/stacks-sites/pinned-original-source.tex, source lines 7566--7566. Navigation pointer invoked at proof line 7566; the counted span contains the pointer only and the cited result is not reproduced inside it, so the pointer is registered as a non-load-bearing navigation stub in the established form for Stacks material. No bibliography citation is involved.\par\medskip
\ifdefined\hypertarget\hypertarget{sourcepointer-3}{}\fi
\noindent\textit{Source reference.} Original lemma 00Y1: sources/stacks-sites/pinned-original-source.tex, source lines 7582--7582. Navigation pointer invoked at proof line 7582; the counted span contains the pointer only and the cited result is not reproduced inside it, so the pointer is registered as a non-load-bearing navigation stub in the established form for Stacks material. No bibliography citation is involved.\par\medskip
\ifdefined\hypertarget\hypertarget{sourcepointer-4}{}\fi
\noindent\textit{Source reference.} Original lemma 03EE: sources/stacks-sites/pinned-original-source.tex, source lines 7590--7590. Navigation pointer invoked at proof line 7590; the counted span contains the pointer only and the cited result is not reproduced inside it, so the pointer is registered as a non-load-bearing navigation stub in the established form for Stacks material. No bibliography citation is involved.\par\medskip
\begin{lemma}
\label{lemma-localize-topos}
Let $\mathcal{C}$ be a site.
Let $\mathcal{F}$ be a sheaf on $\mathcal{C}$.
Then the category $\Sh(\mathcal{C})/\mathcal{F}$
is a topos. There is a canonical morphism of topoi
$$
j_\mathcal{F} :
\Sh(\mathcal{C})/\mathcal{F}
\longrightarrow
\Sh(\mathcal{C})
$$
which is a localization as in
Section \ref{section-localize}
such that
\begin{enumerate}
\item the functor $j_\mathcal{F}^{-1}$ is the functor
$\mathcal{H} \mapsto \mathcal{H} \times \mathcal{F}/\mathcal{F}$, and
\item the functor $j_{\mathcal{F}!}$ is the forgetful
functor $\mathcal{G}/\mathcal{F} \mapsto \mathcal{G}$.
\end{enumerate}
\end{lemma}

\begin{proof}
Apply
Lemma \ref{lemma-topos-good-site}.
This means we may assume $\mathcal{C}$ is a site
with subcanonical topology, and $\mathcal{F} = h_U = h_U^\#$
for some $U \in \Ob(\mathcal{C})$.
Hence the material of
Section \ref{section-localize}
applies. In particular, there is an equivalence
$\Sh(\mathcal{C}/U) = \Sh(\mathcal{C})/h_U^\#$
such that the composition
$$
\Sh(\mathcal{C}/U)
\to
\Sh(\mathcal{C})/h_U^\#
\to \Sh(\mathcal{C})
$$
is equal to $j_{U!}$, see
Lemma \ref{lemma-essential-image-j-shriek}.
Denote
$a : \Sh(\mathcal{C})/h_U^\# \to \Sh(\mathcal{C}/U)$
the inverse functor, so $j_{\mathcal{F}!} = j_{U!} \circ a$,
$j_\mathcal{F}^{-1} = a^{-1} \circ j_U^{-1}$, and
$j_{\mathcal{F}, *} = j_{U, *} \circ a$. The description of
$j_{\mathcal{F}!}$ follows from the above. The description of
$j_\mathcal{F}^{-1}$ follows from
Lemma \ref{lemma-compute-j-shriek-restrict}.
\end{proof}

\end{document}
